\documentclass[11pt]{article}

% Minimal, compilable starter for the thesis/ folder.
% Replace the title, authors, and sections with your own work,
% then compile this file (e.g. `pdflatex main.tex`).
\usepackage{amsmath,amssymb}
\usepackage[margin=1in]{geometry}
\usepackage{graphicx}
\usepackage[colorlinks=true,allcolors=blue]{hyperref}

\title{A Minimal Research Paper}
\author{Your Name\thanks{Affiliation. \texttt{you@example.com}}}
\date{\today}

\begin{document}
\maketitle

\begin{abstract}
This is a minimal \LaTeX{} starter for the \texttt{thesis/} folder. Replace the
title, authors, and sections with your own work, then compile this file with
\texttt{pdflatex main.tex}.
\end{abstract}

\section{Introduction}
Motivate the problem and state your contributions. Cite related work with
\verb|\cite|, for example~\cite{vaswani2017attention}.

\section{Method}
Show that math typesetting works:
\begin{equation}
  \mathcal{L}(\theta)
  = \frac{1}{N} \sum_{i=1}^{N} \ell\bigl(f_{\theta}(x_i),\, y_i\bigr).
  \label{eq:loss}
\end{equation}

\section{Experiments}
Describe the setup, then reference Equation~\eqref{eq:loss} when reporting
results. Add tables and figures as needed.

\section{Conclusion}
Summarize the main takeaway and future work.

\begin{thebibliography}{9}
\bibitem{vaswani2017attention}
A.~Vaswani et al.
\newblock Attention is all you need.
\newblock In \emph{Advances in Neural Information Processing Systems}, 2017.
\end{thebibliography}

\end{document}
